\par\noindent\textbf{(1)} 双光子共振且适当选择相位后，相互作用哈密顿量可写H\_I=−ℏ[Ωp|3><1|+Ωc|3><2|+h.c.]/2（单光子失谐项可置零）。激发态只与亮态Ωp|1>+Ωc|2>耦合，存在与之正交的线性组合。
\par\noindent\textbf{(2)} 暗态|D>=(Ωc|1>−Ωp|2>)/√(|Ωc|²+|Ωp|²)，满足<3|H\_I|D>=0，因此理想绝热跟随时激发态占有接近零，可抑制自发辐射吸收。弱探针时|Ωp|≪|Ωc|，多数人口仍在|1>。
\par\noindent\textbf{(3)} 介质色散k(ω)=n(ω)ω/c，群速v\_g=[dk/dω]\textasciicircum{}\{-1\}=c/[n+ω dn/dω]。EIT透明窗附近探针折射率实部对频率斜率大且为正，产生大的群折射率n\_g=n+ω dn/dω与v\_g≪c；在理想线性响应中群延迟随光学深度增加、随|Ωc|²减小而增大。
\par\noindent\textbf{(4)} 把控制场降到零会使透明窗口收窄，若不满足绝热条件可能恢复吸收，不能简单让v\_g公式预测无限储光。有限基态退相干γ12和激发态自发衰减限制透明深度与延迟带宽积，弱探针近似外也出现非线性饱和。
